\documentclass[11pt]{article}

\usepackage[a4paper,margin=1in]{geometry}
\usepackage{amsmath,amssymb,amsthm,mathtools}
\usepackage{enumitem}
\usepackage{xcolor}
\usepackage{hyperref}

\hypersetup{
  colorlinks=true,
  linkcolor=blue,
  urlcolor=blue,
  citecolor=blue
}

\newtheorem{theorem}{Theorem}
\newtheorem{lemma}[theorem]{Lemma}
\newtheorem{proposition}[theorem]{Proposition}
\theoremstyle{definition}
\newtheorem{hypothesis}{Hypothesis}
\setcounter{hypothesis}{-1}
\newtheorem{definition}{Definition}
\theoremstyle{remark}
\newtheorem{remark}{Remark}

\DeclareMathOperator{\cone}{cone}
\DeclareMathOperator{\relint}{relint}
\DeclareMathOperator{\supp}{supp}
\DeclareMathOperator*{\wdeg}{w\text{-}deg}
\newcommand{\inw}{\operatorname{in}_w}
\newcommand{\R}{\mathbb{R}}

\title{Face-Selection Law at a Simple Polyhedral Vertex}
\author{}
\date{}

\begin{document}

\maketitle

Theorem~V.16 in \texttt{FORMAL\_NEWTON\_TROPICAL.md} selects
$q^\star=\min_j q_j$ over the monomials of the perturbation on the positive
orthant. Section~15 of \texttt{FORMAL\_VERTEX\_THRESHOLD.md} gives the
weighted degree of a single monomial. This note transports both results to a
simple vertex of a general polyhedron, which is what domain three adjudicates.

The transport is a change of coordinates, not a new theorem. What is new is
that the hypotheses it needs are stated explicitly---including several the
first draft left implicit---and that one of them is measured per corpus row
rather than assumed.

\section{Standing assumptions}

$P\subset\R^n$ is a bounded full-dimensional polyhedron; $F$ and $G$ are
continuous on $P$ and real analytic near the vertex under study. Boundedness
is not cosmetic: it is what makes Lemma~\ref{lem:isolation} available, and the
domain's proposal prompt already rejects unbounded systems.

\section{Setup}

Let $v$ be a \emph{simple} vertex with inward edge generators
$u_1,\dots,u_n$, and set
\[
  \Phi(c)=v+\sum_{i=1}^{n}c_iu_i,\qquad c\in\R_{\ge0}^n .
\]
Simplicity makes the $u_i$ linearly independent, so $\Phi$ is a linear
isomorphism from the nonnegative orthant onto the tangent cone $T_vP$; near
$v$ it is a bijection onto a neighbourhood of $v$ in $P$. Put
\[
  D(c)=F(v)-F\bigl(\Phi(c)\bigr),
  \qquad
  R(c)=G\bigl(\Phi(c)\bigr)-G(v),
\]
and study
\[
  M(s)=\sup_{x\in P}\bigl(F(x)+sG(x)\bigr),\qquad s\to0^+ .
\]

\section{Global isolation}

The first draft asserted that $v$ is ``the relevant unperturbed maximiser''
and then worked locally. That is an assumption and needs to be one, because a
perturbation can move the global optimum to a different vertex entirely.

\begin{hypothesis}[isolation]\label{hyp:iso}
$v$ is the unique maximiser of $F$ on $P$.
\end{hypothesis}

\begin{lemma}\label{lem:isolation}
Under Hypothesis~\ref{hyp:iso}, for every neighbourhood $U$ of $v$ there is
$s_U>0$ such that for all $0<s<s_U$ the supremum defining $M(s)$ is attained
in $U$.
\end{lemma}

\begin{proof}
$P\setminus U$ is compact and $F$ is continuous with a unique maximum at $v$,
so
\[
  \eta:=F(v)-\max_{P\setminus U}F>0 .
\]
Let $K=2\sup_P|G|$, finite by compactness. For $x\in P\setminus U$,
\[
  F(x)+sG(x)\;\le\;F(v)-\eta+\tfrac{sK}{2}\;<\;F(v)+sG(v)
\]
as soon as $sK<\eta$. The value at $v$ already exceeds everything outside
$U$, so the supremum is attained in $U$; take $s_U=\eta/K$.
\end{proof}

Everything below is therefore about the local problem in edge coordinates, and
Lemma~\ref{lem:isolation} is what licenses that reduction. Without
Hypothesis~\ref{hyp:iso}---at a tie between two vertices, say---the conclusion
can fail outright, which is part of what the domain's
\texttt{linear\_max\_at\_vertex} rule exists to exercise.

\section{Weights and initial forms}

Fix the weight vector $w=(1/\beta_1,\dots,1/\beta_n)$, the $\beta_i$ pinned by
Hypothesis~\ref{hyp:base} below. For a monomial $c^{\alpha}$,
\[
  \wdeg\bigl(c^{\alpha}\bigr)=\langle w,\alpha\rangle
  =\sum_{i=1}^{n}\frac{\alpha_i}{\beta_i}.
\]

\begin{definition}[initial form]
For $H=\sum_j\gamma_jc^{\alpha_j}\not\equiv0$, the \emph{initial form}
$\inw(H)$ is the sum of those terms of least $w$-degree, and $\wdeg(H)$ is
that least degree.
\end{definition}

The anisotropic dilation
\[
  \delta_\tau(c)=\bigl(\tau^{1/\beta_1}c_1,\dots,\tau^{1/\beta_n}c_n\bigr)
\]
sends $c^{\alpha}$ to $\tau^{\langle w,\alpha\rangle}c^{\alpha}$, so $\inw(H)$
is exactly what survives in $\tau^{-\wdeg(H)}H(\delta_\tau c)$ as
$\tau\to0^+$. Both $D_0$ and $W_S$ below are initial forms in this sense,
not descriptions of one.

\section{Hypotheses, with uniform remainders}

\begin{hypothesis}[positivity]\label{hyp:pos}
$D(c)\ge0$ near $c=0$, with equality only at $c=0$.
\end{hypothesis}

\begin{hypothesis}[weighted principal part, uniform]\label{hyp:base}
There are $A_i>0$ and $\beta_i>1$ such that, with
$D_0(c)=\sum_iA_ic_i^{\beta_i}=\inw(D)$,
\[
  D(\delta_\tau z)=\tau\bigl(D_0(z)+e_D(\tau,z)\bigr),
  \qquad
  \sup_{z\in K}\lvert e_D(\tau,z)\rvert\longrightarrow0
  \quad(\tau\to0^+)
\]
for every compact $K$ in the orthant.
\end{hypothesis}

\begin{hypothesis}[polynomial perturbation, uniform facewise]\label{hyp:pert}
$R(c)=\sum_j\gamma_jc^{\alpha_j}$ is a finite sum, and for every face $S$ on
which $R$ does not vanish identically,
\[
  R|_{C_S}(\delta_\tau z)=\tau^{q_S}\bigl(W_S(z)+e_R(\tau,z)\bigr),
  \qquad
  \sup_{z\in K}\lvert e_R(\tau,z)\rvert\longrightarrow0
\]
for every compact $K\subset C_S$, with $q_S$ and $W_S$ as in
Section~\ref{sec:faces}.
\end{hypothesis}

Uniformity is the point. The first draft wrote $o(\tau)$ pointwise in $z$ and
then took a supremum over $z$, which does not follow: pointwise remainders
permit the error to blow up along a sequence of $z$ as $\tau\to0$. With $R$ a
finite polynomial and $D$ analytic the uniform version holds automatically on
compacta, so nothing is lost---but it must be said, because the balance in
Section~\ref{sec:balance} exchanges a limit with a supremum.

\begin{remark}
Hypothesis~\ref{hyp:base} is the one the transport smuggles in, and the only
one not implied by the setting: an arbitrary $F$ written in edge coordinates
may carry cross terms of lower weight. The quantity
\texttt{base\_homogeneity} in the corpus measures the exponent of
$D(\delta_\tau z)$; the value $1$ means Hypothesis~\ref{hyp:base} holds for
that row.
\end{remark}

\section{Faces, and admissibility defined from the data alone}
\label{sec:faces}

For nonempty $S\subseteq\{1,\dots,n\}$ let
\[
  C_S=\{c\ge0:\ c_i=0\ \text{for}\ i\notin S\},
\]
which $\Phi$ carries to $v+\cone\{u_i:i\in S\}$, the corresponding face of the
tangent cone. Monomial $j$ has support $S_j=\{i:\alpha_{ij}>0\}$, and
$c^{\alpha_j}$ is nonzero on $\relint(C_S)$ exactly when $S_j\subseteq S$.
Monomial support and tangent-cone face are the same data in edge coordinates.

Every $c\ge0$ lies in $\relint(C_S)$ for exactly one $S$, namely $\supp(c)$,
so the faces partition the cone and
\begin{equation}\label{eq:partition}
  \sup_{c\ge0}\;=\;\max_{S}\ \Bigl(\ \sup_{\relint(C_S)}\ \Bigr).
\end{equation}

\begin{definition}[from the data $(F,G,T_vP)$ alone]\label{def:adm}
For each face $S$:
\begin{itemize}[nosep]
  \item $S$ is \textbf{active} if $R|_{C_S}\not\equiv0$, i.e.\ some
        $S_j\subseteq S$. For active $S$ put $q_S=\wdeg(R|_{C_S})$ and
        $W_S=\inw(R|_{C_S})$; equivalently
        $q_S=\min\{q_j:\ S_j\subseteq S,\ \gamma_j\ne0\}$ with
        $q_j=\langle w,\alpha_j\rangle$.
  \item $S$ is \textbf{positive} if $W_S(z)>0$ for some $z\in\relint(C_S)$.
  \item $S$ is \textbf{subcritical} if $q_S<1$.
\end{itemize}
$S$ is \textbf{admissible} when it is active, positive and subcritical.
\end{definition}

None of the three properties refers to $M(s)$, to $\gamma$, or to any
conclusion of this note; all are computable from the exponents, the
coefficients and the cone. That is what makes Section~\ref{sec:selection}
non-circular. The first draft folded ``the relevant branch has positive
effective perturbation and $q^\star<1$'' into a hypothesis and then used
admissibility in the conclusion, which is close to assuming what is to be
proved.

\begin{hypothesis}[no leading cancellation]\label{hyp:nocancel}
For each admissible $S$, $W_S$ does not vanish identically on $\relint(C_S)$.
\end{hypothesis}

With $W_S$ a nonzero polynomial this is automatic; it is stated because the
initial form of a \emph{sum} over a face can cancel even when no individual
monomial does.

\section{Inactive and non-positive faces contribute nothing}

The first draft excluded such faces by fiat. Under \eqref{eq:partition} they
must be shown harmless, because the supremum ranges over all of them.

\begin{lemma}\label{lem:inadmissible}
Let $S$ be a face that is not admissible. Then $S$ yields no positive
improvement for all small $s$, and in particular never determines the leading
order.
\end{lemma}

\begin{proof}
Three cases.

\emph{Inactive.} $R|_{C_S}=0$, so the objective on $C_S$ is $-D(c)\le0$ by
Hypothesis~\ref{hyp:pos}, with supremum $0$ approached only as $c\to0$. It
contributes nothing positive at all. This is the tilted simplex's vertical
edge: feasible, but perturbatively silent.

\emph{Active, not positive.} $W_S\le0$ throughout $\relint(C_S)$. Writing
$c=\delta_\tau z$ with $z$ on a compact cross-section,
Hypothesis~\ref{hyp:pert} gives
$R|_{C_S}(\delta_\tau z)=\tau^{q_S}(W_S(z)+e_R)\le0$ for small $\tau$ by
uniformity, whence the objective is $\le-\tau(D_0(z)+e_D)\le0$.

\emph{Active, positive, not subcritical.} $q_S\ge1$. The leading balance of
Section~\ref{sec:balance} is $-\tau A+s\tau^{q_S}B$ with $A,B>0$. For
$q_S>1$, $\tau^{q_S}=o(\tau)$ and the expression is negative for all small
$\tau>0$ once $s$ is small. For $q_S=1$ it equals $\tau(sB-A)$, negative once
$s<A/B$. Either way the face yields no positive improvement at small $s$.
\end{proof}

Lemma~\ref{lem:inadmissible} turns \eqref{eq:partition} into a maximum over
admissible faces only. It is the step that makes ``inadmissible'' mean
\emph{contributes nothing} rather than \emph{excluded by hand}.

\section{The facewise balance}\label{sec:balance}

Fix an admissible $S$ and a compact cross-section $Z$ of $\relint(C_S)$---for
instance $\{z\in C_S:D_0(z)=1\}$, compact by Hypothesis~\ref{hyp:pos}. Write
$c=\delta_\tau z$. By Hypotheses~\ref{hyp:base} and~\ref{hyp:pert}, uniformly
for $z\in Z$,
\[
  J_s(\tau,z)=-D(\delta_\tau z)+sR(\delta_\tau z)
  =-\tau\bigl(D_0(z)+e_D\bigr)+s\tau^{q_S}\bigl(W_S(z)+e_R\bigr).
\]
With $A=D_0(z)>0$, $B=W_S(z)>0$ and $k=q_S\in(0,1)$, the leading expression
$-\tau A+s\tau^{k}B$ is stationary at
\[
  \tau_\ast=\Bigl(\frac{skB}{A}\Bigr)^{\!1/(1-k)},
\]
and substituting back gives
\[
  J=A\,\tau_\ast\,\frac{1-k}{k}\;>\;0 .
\]
Since $\tau_\ast=\Theta\bigl(s^{1/(1-k)}\bigr)$, so is $J$. Uniformity of
$e_D$ and $e_R$ on $Z$ is what lets the supremum over $z$ pass inside the
limit, giving matching upper and lower constants and hence
\[
  \sup_{c\in\relint(C_S)}J_s=\Theta\bigl(s^{1/(1-q_S)}\bigr),
  \qquad
  \gamma_S=\frac{1}{1-q_S}.
\]

\begin{remark}[checked against measurement]
For the tilted simplex face $c_2=0$ ($D_0=z^4$, $R=z$, $k=1/4$) the
coefficient is $\tfrac34(\tfrac14)^{1/3}=0.472470$, and the adjudicator
measures $0.472470\,s^{4/3}$ to six decimals at
$s=10^{-2},\,2.5\times10^{-3},\,6.25\times10^{-4}$.
\end{remark}

\section{Selection, qualified}\label{sec:selection}

\begin{theorem}\label{thm:selection}
Assume Hypotheses~\ref{hyp:iso}--\ref{hyp:nocancel} and that at least one face
is admissible. Put $q^\star=\min\{q_S:\ S\ \text{admissible}\}$. Then
\[
  M(s)-F(v)-sG(v)=\Theta\bigl(s^{1/(1-q^\star)}\bigr),
  \qquad
  \gamma=\frac{1}{1-q^\star}.
\]
\end{theorem}

\begin{proof}
Lemma~\ref{lem:isolation} localises the supremum to a neighbourhood of $v$,
where $\Phi$ is a bijection onto the cone. By \eqref{eq:partition} and
Lemma~\ref{lem:inadmissible} the supremum is the maximum over admissible faces
of $\Theta(s^{1/(1-q_S)})$ from Section~\ref{sec:balance}. Since $1/(1-q)$ is
increasing on $(0,1)$ and $s<1$, the largest such term is the one with the
smallest exponent, i.e.\ the smallest $q_S$.
\end{proof}

\begin{proposition}[relation to V.16's monomial minimum]\label{prop:v16}
$q^\star\ge\min_j q_j$, with equality precisely when some face attaining
$\min_j q_j$ is admissible.
\end{proposition}

\begin{proof}
$q_S$ is a minimum over the monomials supported in $S$, so the minimum of
$q_S$ over \emph{all} nonempty faces is attained at the full face
$S=\{1,\dots,n\}$ and equals $\min_j q_j$. Restricting the minimum to
admissible faces can only raise it, which gives the inequality; equality holds
exactly when the restricted minimum still reaches $\min_j q_j$, i.e.\ when
some face attaining it is admissible.
\end{proof}

V.16's statement is therefore recovered under a sufficient condition---for
example all $\gamma_j>0$, whence every active face is positive and no initial
form cancels, together with $\min_j q_j<1$---and \emph{not} in general.

\begin{remark}
The qualification is not pedantic. \texttt{3d\_shear\_orders\_2\_4\_6} in the
corpus has an inadmissible face carrying $q=1.3804$ alongside admissible ones
at $q=0.5$. There the inadmissible face is the larger value, so nothing is
lost; but that ordering is not guaranteed, and an implementation taking
$\min_jq_j$ blindly would in general select a branch that contributes nothing.
The adjudicator computes the face minimum, which is the quantity
Theorem~\ref{thm:selection} is about.
\end{remark}

\section{Worked case: the tilted simplex}

Let $P=\{x_0\ge0,\ x_1\ge0,\ x_0+x_1\le1\}$ and $v=(0,1)$. The active
constraints at $v$ are $x_0\ge0$ and $x_0+x_1\le1$, so
\[
  T_vP=\{d\in\R^2:\ d_0\ge0,\ d_0+d_1\le0\},
\]
with extreme rays $u_1=(1,-1)$ and $u_2=(0,-1)$, and
$x=(c_1,\,1-c_1-c_2)$. For $F=-\bigl((x_0+x_1-1)^2+x_0^4\bigr)$ we have
$x_0+x_1-1=-c_2$ and $x_0=c_1$, hence
\[
  D(c)=c_2^2+c_1^4,\qquad \beta=(4,2),\qquad w=(\tfrac14,\tfrac12).
\]
With $G=x_0$ we get $R(c)=c_1$, of support $\{1\}$. Face $\{1\}$ is active,
positive and subcritical with $q=1/4$. Face $\{2\}$ is \textbf{inactive}: $R$
vanishes on it identically, so by Lemma~\ref{lem:inadmissible} its $\beta=2$
never enters---the quadratic term is the leading behaviour of $F$ along that
edge and is nonetheless irrelevant, which is the point the first draft got
wrong by claiming the quadratic term was constant on the feasible cone. Hence
$q^\star=1/4$ and $\gamma=4/3$.

The ambient axis $e_0=(1,0)$ has $e_{0,0}+e_{0,1}=1>0$, so it is not in
$T_vP$ at all; its quadratic decay is not a branch of the constrained problem.
That is the \texttt{ambient\_exponent\_law} counterexample, stated exactly.

\section{Scope}

Conditional statement: for a \emph{simple} vertex, under
Hypotheses~\ref{hyp:iso}--\ref{hyp:nocancel}, the asymptotic exponent is
$\gamma=1/(1-q^\star)$ with $q^\star$ the minimum weighted degree over
\emph{admissible} faces.

Not established here:

\begin{itemize}
  \item \textbf{Hypothesis~\ref{hyp:base} for arbitrary $F$.} In edge
        coordinates a base may carry cross terms; one of lower weight makes
        the drop $\Theta(\tau^{c})$ with $c<1$, and the per-edge $\beta_i$
        stop describing the cone's interior. Constructed violations at
        $c=0.75$ and $c=0.5$ still predict correctly---a positive cross term
        only steepens the interior and drives the optimum onto a face---but
        agreeing is not being licensed. This is why
        \texttt{base\_homogeneity} is recorded per row and never gated on.
  \item \textbf{Non-simplicial vertices.} With more than $n$ active
        constraints $\Phi$ is not an isomorphism onto the orthant and the
        argument fails at Section~2. The vertex probe returns the vertex value
        there rather than answering.
  \item \textbf{Hypothesis~\ref{hyp:iso} in the corpus.} Uniqueness of the
        unperturbed maximiser is measured rather than assumed, but only by a
        finite probe. \texttt{rival\_margin} records per row the amount by
        which the vertex outranks the best competing maximum found outside a
        ball of radius $0.25$; it can exhibit a rival, it cannot certify that
        none exists, so it is recorded and never gated on, exactly as
        \texttt{base\_homogeneity} is. The count is not zero: the generated
        report \texttt{docs/POLYHEDRA.md} lists the rows that come back at or
        below zero, and some of them are rows the edge law is otherwise
        counted as confirming. There the maximiser is a whole face,
        Lemma~\ref{lem:isolation} does not hold, and the agreement between
        predicted and measured exponent is not evidence for the law.
\end{itemize}

None of the three is a gap in the coordinate transport; all three are limits
on where it applies.


\section{What the framework answers}

This section is an index, not an argument. Each entry names a question the
framework is asked in practice and points at the result above that settles it.
Nothing here is new; the purpose is to make the correspondence explicit for a
reader who arrives with the question rather than with the proof.

\paragraph{Edge generators and simplicity.}
At a simple vertex the inward edge generators $u_1,\dots,u_n$ are linearly
independent by definition, so $\Phi$ is a linear isomorphism from the
nonnegative orthant onto $T_vP$ and a bijection onto a neighbourhood of $v$ in
$P$. In the tilted simplex $u_1=(1,-1)$ and $u_2=(0,-1)$, and simplicity is
$\det\left(\begin{smallmatrix}1&0\-1&-1\end{smallmatrix}\right)=-1\ne0$.

\paragraph{Tangent cone and the monomial--face dictionary.}
Orthant faces $C_S$ map to the faces $v+\cone\{u_i:i\in S\}$ of the tangent
cone, and $c^{\alpha_j}$ survives on $\relint(C_S)$ exactly when
$S_j\subseteq S$. Monomial support and tangent-cone face are the same
combinatorial datum in edge coordinates; the minimal supporting face of
monomial $j$ is $C_{S_j}$.

\paragraph{Weighted principal part.}
Hypothesis~\ref{hyp:base}, with the remainder uniform on compacta. Under
$\delta_\tau$ one has the exact homogeneity $D_0(\delta_\tau c)=\tau D_0(c)$.
The corpus quantity \texttt{base\_homogeneity} measures the resulting scaling
exponent; the value $1$ licenses the hypothesis for that row. In the tilted
simplex it holds exactly, with $D_0(c)=c_1^4+c_2^2$. The principal part is a
weighted-homogeneous power sum, not a linear form.

\paragraph{Weighted degree of a monomial.}
Each monomial transforms under $\delta_\tau$ by $\tau^{q_j}$ with
$q_j=\langle w,\alpha_j\rangle=\sum_i\alpha_{ij}/\beta_i$: the
Newton--tropical valuation with respect to the base weights $w_i=1/\beta_i$.
The value belongs to one monomial, hence to one minimal supporting face. It is
\emph{not} obtained by summing independently measured ray exponents, which is
the error this note exists to correct.

\paragraph{Restriction, initial forms and cancellation.}
On $C_S$ only monomials with $S_j\subseteq S$ survive; $q_S$ is the least of
their weighted degrees, and $W_S=\inw(R|_{C_S})$ is an explicit homogeneous
polynomial. If that lowest-weight sum vanishes identically one must pass to
the next layer, which is what Hypothesis~\ref{hyp:nocancel} excludes on faces
later called admissible.

\paragraph{Admissibility, and why it is non-circular.}
Definition~\ref{def:adm}: active, positive, subcritical. All three refer only
to $(D_0,R)$ and the combinatorial face, never to $M(s)$ or to $\gamma$.
Lemma~\ref{lem:inadmissible} then shows that a face failing any of them
contributes nothing positive at small $s$, so ``inadmissible'' means
\emph{contributes nothing} rather than \emph{excluded by hand}.

\paragraph{Selection and the qualified minimum.}
Theorem~\ref{thm:selection} takes $q^\star$ over admissible faces, and
Proposition~\ref{prop:v16} relates it to V.16's monomial minimum:
$q^\star\ge\min_jq_j$, with equality exactly when some face attaining
$\min_jq_j$ is admissible. If a formally smaller $q_j$ belongs to an
inadmissible face it is simply omitted, and the asymptotics are governed by
the next-smallest admissible value.

\paragraph{Predicted exponent, stationary scale, exact coefficient.}
The facewise balance of Section~\ref{sec:balance} gives
$\gamma_S=1/(1-q_S)$, and Lemma~\ref{lem:isolation} promotes the dominant face
to the global statement. For the tilted simplex $\gamma=4/3$, and solving the
one-dimensional problem on the active face gives the exact identity
$M(s)=3\cdot4^{-4/3}s^{4/3}$. The anisotropic stationary scale is
$\tau_\ast\asymp s^{4/3}$ while the physical edge coordinate scales as
$c_1^\ast\asymp s^{1/3}$; the two are related by $c_1=\tau^{1/4}z_1$.

\paragraph{Inactive directions that look plausible.}
Two are correctly excluded: the vertical edge, geometrically feasible but with
$R\equiv0$ on it, and the ambient axis $e_0=(1,0)$, which leaves $P$
immediately and so belongs to no face of $T_vP$. Ambient exponents computed
along infeasible rays have no standing in the constrained problem.

\section{Three principles}

The results above organise into three interlocking statements. They are
restatements, not additions: each is a reading of a lemma already proved.

\subsection*{1. The boundary controls the singularity}

Near a simple vertex the singular behaviour of $M(s)-F(v)-sG(v)$ is determined
by the faces of the tangent cone, not by the relative interior of the full
cone.

After transport to edge coordinates the base drop has principal part
$D_0(c)=\sum_iA_ic_i^{\beta_i}$ with $\beta_i>1$, so any interior direction
(all $c_i>0$) incurs a strictly positive combination of these costs. A
perturbation monomial of weighted degree $q$ can produce a positive
improvement of order $s^{1/(1-q)}$ only when $q<1$, and the balance realising
it is the elementary competition between a term of order $\tau$ and one of
order $s\tau^{q}$.

The same competition can be run on an open subset of the orthant, but its
critical point can always be projected onto a face by setting coordinates to
zero. Every coordinate set to zero removes a non-negative base cost while, by
the support rule, never introducing a leading perturbation term of lower
weight, so the projected point improves by at least as much. Consequently the
maximiser of the local model lies on a face---which is what
\eqref{eq:partition} together with Lemma~\ref{lem:inadmissible} says
formally.

In the tilted simplex this is visible by elementary calculus: the objective
$-a^4-b^2+sa$ is strictly decreasing in $b$ for $b>0$, so the maximiser is
forced onto the slanted edge $b=0$. The interior of the cone is not irrelevant
to the original optimisation; it is silent as a \emph{distinct leading
mechanism}, because its contribution is matched or improved by an admissible
boundary face.

\subsection*{2. The lowest weight dominates}

Among admissible faces the smallest $q_S$ produces the dominant gap, of order
$s^{1/(1-q_S)}$; higher-weight faces contribute only lower-order terms.

On an admissible face the balance $-\tau A+s\tau^{q_S}B$ with $A,B>0$ has
stationary value of exact order $s^{1/(1-q_S)}$, and $q\mapsto1/(1-q)$ is
strictly increasing on $(0,1)$. A smaller $q$ therefore gives a strictly
larger exponent and, for small $s$, a strictly larger improvement. Comparing
faces gives $q^\star=\min_{S\ \mathrm{admissible}}q_S$, every larger $q_S$
contributing a remainder $o(s^{1/(1-q^\star)})$.

The ordering recurs one level down: $q_S$ is itself the minimum over the
monomials surviving on $S$. So the lowest weight, once restricted to
admissible faces, is the only quantity entering $\gamma=1/(1-q^\star)$. There
is no averaging and no max-weight rule; the Newton--tropical selection is a
pure minimum.

\subsection*{3. The feasible face selects the exponent}

Only faces inside the tangent cone that satisfy Definition~\ref{def:adm}
compete. Geometry (feasibility) and algebra (positive initial form,
$0<q_S<1$) together select the exponent; ambient or inactive directions are
invisible to the asymptotics.

A direction leaving the polyhedron---the ambient axis $e_0=(1,0)$ at the
tilted-simplex vertex---belongs to no face of $T_vP$ and never enters the
minimum. A geometrically feasible face on which the perturbation vanishes
identically---the vertical edge $C_{\{2\}}$---is inactive and is excluded on
the same terms. After these exclusions a finite admissible set remains, and
the exponent is read from it alone. The numerical value of an ambient or
inactive degree, however suggestive, cannot alter the prediction: the exponent
is a combinatorial-algebraic invariant of the admissible faces of the tangent
cone. Feasibility is not an afterthought, it is part of the selection rule.

\subsection*{Unified statement}

Taken together: the singular asymptotics of a linearly perturbed maximisation
problem near a simple polyhedral vertex are determined by a minimum of
Newton--tropical weights, taken only over those faces of the tangent cone that
are both geometrically feasible and algebraically capable of producing a
positive improvement of order less than one. The boundary carries the
singularity, the lowest admissible weight dominates, and the feasible face
selects the exponent.

That turns a technical asymptotic calculation into a selection principle with
three layers:
\begin{enumerate}[nosep]
  \item \emph{Localisation}: the tangent cone replaces the global geometry
        (Lemma~\ref{lem:isolation}).
  \item \emph{Selection}: admissible faces are ranked by weighted degree
        (Definition~\ref{def:adm}, Lemma~\ref{lem:inadmissible}).
  \item \emph{Scaling}: the winning degree becomes the response exponent
        (Section~\ref{sec:balance}, Theorem~\ref{thm:selection}).
\end{enumerate}

The boundary principle is not merely the observation that the optimiser
happens to lie on a boundary. It says the boundary is the space in which the
asymptotic competition is \emph{selected}. The face lattice of $T_vP$ acts as
a discrete set of asymptotic channels and the variational problem chooses
among them. A face wins not because it has fewer variables but because
removing directions can lower the base cost without sacrificing the relevant
perturbative gain, which gives the principle a variational rather than merely
geometric meaning.

The second principle has a consequence worth isolating: the exponent is
insensitive to much that affects the coefficient. Once the winning admissible
degree is known, $\gamma=1/(1-q^\star)$ is fixed, and distinct functions,
coefficients or polyhedra sharing that degree share the exponent. The reduced
optimisation on the winning face determines only the prefactor.

The third is conceptually decisive. Newton data cannot be read in the ambient
coordinate space alone; the relevant Newton diagram is filtered through the
tangent-cone face lattice. An ambient monomial may look dominant
algebraically and contribute nothing, because its direction is infeasible or
its restriction to the feasible face vanishes. The slogan is that
\emph{the Newton minimum is taken after geometric restriction, not before}.
That is what prevents incorrect predictions from ambient scaling arguments,
and it is the content of the \texttt{ambient\_exponent\_law} counterexample.

The hierarchy implicit in all this is
\[
  \text{ambient monomials}\longrightarrow\text{face restrictions}
  \longrightarrow\text{admissible weights}\longrightarrow q^\star
  \longrightarrow\gamma,
\]
each arrow discarding information that cannot influence the leading
singularity. What survives is a small invariant: the minimum admissible face
weight. The law is best read as a local asymptotic analogue of an active-set
rule,
\[
  \text{active constraint geometry}+\text{weighted perturbation order}
  \;\Longrightarrow\;\text{singular response exponent},
\]
which explains why the power is encoded by the interaction between the Newton
weights and the feasible face structure rather than merely computing it in one
example.

\section{Implications}

\begin{enumerate}[wide=0pt]

\item \textbf{Exponents become geometric invariants.}
The response exponent is not primarily set by the numerical coefficients of
$D$ and $R$ but by the weighted degree of the winning admissible face,
$\gamma=1/(1-q^\star)$. Two optimisation problems sharing $q^\star$ belong to
the same asymptotic universality class. Coefficients move the prefactor; the
exponent is stable under many perturbations of the data.

\item \textbf{Singular response is an active-constraint phenomenon.}
A fractional-power law signals that the optimiser is asymptotically
interacting with the boundary of the feasible cone, so the exponent records
which constraints become active as well as the analytic order of the
perturbation. An observed $s^{4/3}$ law points to an effective weighted degree
$q^\star=1/4$; a different exponent indicates a different active face or
scaling regime.

\item \textbf{Newton theory must be constrained by geometry.}
The relevant Newton--tropical object is the perturbation \emph{after}
restriction to feasible faces. The correct order is: restrict to a feasible
face, compute the surviving weights, then take the minimum. An ambient
monomial can be algebraically dominant and contribute nothing because it
points outside the tangent cone or vanishes on the active face.

\item \textbf{Boundary stratification gives an algorithm.}
For a simple polyhedral vertex the analysis is a finite combinatorial
procedure:
\begin{enumerate}[nosep,label=(\alph*)]
  \item enumerate the faces of $T_vP$;
  \item restrict the perturbation to each;
  \item discard faces where the initial form vanishes or cannot be positive,
        and those with $q_S\ge1$;
  \item compute $q_S$ on the remainder;
  \item select $q^\star=\min q_S$;
  \item solve the reduced optimisation on the winning face for the
        coefficient.
\end{enumerate}
A local variational problem becomes a face selection followed by a
lower-dimensional calculation. This is what the adjudicator does.

\item \textbf{Adding constraints can change the universality class.}
A new constraint can remove the previously winning face; the minimum is then
taken over a different admissible set and the exponent jumps from
$1/(1-q_1)$ to $1/(1-q_2)$. Response behaviour can therefore change sharply
under geometric deformation even when $F$ and $G$ vary smoothly, the
transition being caused by a change in the active face lattice.

\item \textbf{Perturbations can be classified by relevance.}\label{imp:classes}
The admissible weights sort perturbation terms into asymptotic classes:
$0<q<1$ relevant, of fractional-power type; $q=1$ critical, requiring a
different balance; $q>1$ subleading for this mechanism; vanishing or
sign-indefinite restrictions geometrically inactive. The hierarchy is
analogous to scaling classifications in singular perturbation theory.

\item \textbf{The leading coefficient carries finer geometric information.}
The exponent sees only $q^\star$. The coefficient remembers the shape of the
winning face, the principal form $D_0$ and the initial form $W_S$:
\[
  M(s)-F(v)-sG(v)\sim C_S\,s^{1/(1-q_S)},
\]
with $C_S$ from the reduced optimisation. The theory separates universal
information ($q^\star$) from model-specific information ($C_S$).

\item \textbf{The law is invertible in principle.}
Read backwards, a measured fractional exponent $\gamma$ gives the effective
perturbative weight $q^\star=1-1/\gamma$. Observed scaling can therefore
identify the weighted order of the active perturbation and constrain which
face controls the response.

\item \textbf{It predicts which perturbations matter.}
For $R=\sum_\alpha R_\alpha$, only the lowest-weight term surviving positively
on an admissible face can control the leading response. Higher-weight terms
may alter the prefactor or later corrections but cannot change the first
singular exponent unless the leading term is removed by geometry or
cancellation.

\item \textbf{The phenomenon is locally determined.}
Only the tangent cone, the weighted principal part and the restricted
perturbation enter, so distant parts of the polyhedron do not affect the
leading singularity---a strong locality principle taking the global
optimisation problem to local vertex data. Polyhedral singularities therefore
possess a finite, discrete asymptotic structure: their leading response is a
competition among feasible faces, and the winning face is the channel through
which the perturbation is amplified.

\end{enumerate}

\section{The law as a portable principle}

The face-selection law is not a calculation for one example. Within the
hypotheses of Theorem~\ref{thm:selection} it is a portable principle for
singular asymptotics at simple polyhedral vertices---one that

\begin{itemize}
  \item \textbf{organises the theory} --- localisation to the tangent cone,
        selection among admissible faces, and conversion of the winning weight
        into a response exponent form a single hierarchy;

  \item \textbf{explains the mechanism} --- the boundary controls the
        singularity because interior directions project onto faces without
        losing leading improvement; the lowest admissible weight dominates
        because $q\mapsto1/(1-q)$ is increasing; the feasible face selects the
        exponent because only algebraically and geometrically admissible
        channels enter the minimum;

  \item \textbf{predicts the exponent} --- once the admissible face weights
        are known, $\gamma=1/(1-q^\star)$ with
        $q^\star=\min\{q_S:\ S\ \text{admissible}\}$;

  \item \textbf{filters irrelevant directions} --- ambient axes leaving the
        polyhedron, faces on which the perturbation vanishes, and weights
        $q_S\ge1$ are excluded by Definition~\ref{def:adm} before any
        comparison is made, and Lemma~\ref{lem:inadmissible} shows the
        exclusion costs nothing;

  \item \textbf{classifies perturbations} --- relevant, critical, subleading
        and inactive, as in implication~\ref{imp:classes};

  \item \textbf{reveals active constraints} --- a measured fractional exponent
        identifies both the weighted order of the active perturbation,
        $q^\star=1-1/\gamma$, and the face on which it is realised;

  \item \textbf{generalises cleanly} --- the three principles apply at any
        simple vertex once edge coordinates, a weighted principal part and a
        polynomial perturbation are given; only the combinatorics of the face
        lattice change;

  \item \textbf{is checked, within a stated reach} --- the tilted simplex
        gives the exact identity $M(s)=3\cdot4^{-4/3}s^{4/3}$, the edge law
        carries no counterexample in the corpus while the ambient rule carries
        two, and inactive faces are discarded correctly. What the corpus does
        \emph{not} yet settle is the selection clause itself. Only a row whose
        admissible faces disagree can separate the minimum from the maximum,
        and only one whose rival degree is also below $1$ lets measurement
        choose between two finite numbers rather than merely reject a
        divergence. The adjudicator records that per row as
        \texttt{selection\_discriminates}, alongside
        \texttt{hypotheses\_licensed} and \texttt{rival\_margin}, and
        \texttt{docs/POLYHEDRA.md} reports how many rows satisfy all three at
        once. That count is the number to move.
\end{itemize}

In short: the law converts a local variational problem near a simple
polyhedral vertex into a finite face-selection procedure whose output is a
geometric invariant---the minimum admissible weighted degree---together with a
model-specific prefactor from the reduced optimisation on the winning face.
That is the content of the three principles and of the ten implications that
follow from them.
\end{document}
